% 第七章 总结与展望
\section{总结与展望}

\subsection{工作总结}

本文提出并实现了数字遗产管家系统，一个基于门限密码学与同态承诺的可验证数字资产继承协议。主要工作总结如下：

\begin{enumerate}[label=(\arabic*)]
    \item 设计并实现了双多项式 Shamir 份额结构，在标准 Shamir 方案基础上引入盲因子多项式，使得份额值通过 Pedersen 承诺可在公开存储的同时具备可验证性。
    \item 提出了基于零和多项式的同态份额刷新协议，在 server 完全不存储主密钥的前提下实现安全的份额轮换，填补了现有方案在零信任密钥管理方面的空白。
    \item 设计了计算不可区分的胁迫码机制，为数字资产继承领域提供了首创的抗胁迫方案。
    \item 实现了完整的 Web 系统，覆盖"创建—分发—验证—刷新—恢复"全生命周期，并以 MIT 许可证开源。
\end{enumerate}

\subsection{未来工作}

本系统在未来可以从以下方向进一步扩展：

\textbf{分布式密钥生成（DKG）}。当前方案在创建阶段需要生成完整主密钥（尽管随后销毁），存在理论上创建者单方面获取密钥的瞬时风险。引入 Feldman DKG 协议 \cite{feldman1987practical}，由 $t$ 位监护人通过交互式协议共同生成主密钥，消除对任何单一实体的信任。

\textbf{门限签名与链上资产转移}。当前资产恢复流程为"解密 $\rightarrow$ 发送"，对于区块链资产存在二次转移的成本和安全性问题。引入 ECDSA 门限签名协议（如 GG18/GG20 \cite{gennaro2016threshold}），实现 $t$ 位监护人直接签名链上交易，跳过密钥恢复环节。

\textbf{零知识证明验证}。当前的 Pedersen 承诺同态验证需要公开所有 $t$ 个份额的承诺并进行 Lagrange 系数求和。引入 zk-SNARK/STARK 技术可提供更优的隐私保护和验证效率 \cite{ben2020scalable}。

\textbf{形式化验证}。使用 ProVerif \cite{blanchet2001efficient} 或 Tamarin \cite{meier2013tamarin} 协议验证器对系统的密码学协议进行形式化建模和安全性证明，确保协议的正确性。

\textbf{抗量子密码学迁移}。当前的 secp256k1 曲线可被 Shor 算法在量子计算机上高效破解。随着 NIST 后量子密码标准化进程的推进，将系统迁移到 CRYSTALS-Dilithium \cite{lyubashevsky2022crystals} 等格基密码方案是长期演进方向。

\textbf{安全工程强化}。包括集成硬件安全模块（HSM）、实现前向安全性、完善侧信道攻击防护和通过第三方安全审计。


% 占位说明：如无需独立展望章节，可合并至本章结尾
